% Part: incompleteness
% Chapter: theories-computability
% Section: extensions-of-q-not-decidable

\documentclass[../../../include/open-logic-section]{subfiles}

\begin{document}

\olfileid{inc}{tcp}{cqn}

\olsection{$\Th{Q}$ चे सुसंगत विस्तार अनिर्णेय असतात}

\begin{explain}
कोणत्याही सूत्र~$!A$ साठी $!A$ आणि $\lnot !A$ ही दोन्ही सूत्रे
सिद्ध होत नसतील, तर उपपत्ती \emph{सुसंगत} असते, हे लक्षात ठेवा.
व्याघातापासून कोणतेही वाक्य निष्पन्न होत असल्यामुळे विसंगत उपपत्ती
क्षुल्लक असते: तिच्यात प्रत्येक वाक्य सिद्ध होते.
उपपत्ती $\omega$-सुसंगत असेल, तर ती सुसंगतही असते, हे स्पष्ट आहे.
पण सुसंगतता ही अधिक दुर्बल अट आहे (म्हणजे सुसंगत असलेल्या, पण
$\omega$-सुसंगत नसलेल्या उपपत्ती आहेत). त्यामुळे
\olref[oqn]{thm:oconsis-q} मधील गृहीतक साध्या सुसंगततेपर्यंत
शिथिल करून अधिक सबल प्रमेय मिळवता येते.
\end{explain}

\begin{lem}
``सार्वत्रिक संगणनक्षम संबंध'' अस्तित्वात नाही. म्हणजे पुढील गुणधर्म
असलेला कोणताही द्विस्थानी संगणनक्षम संबंध $R(x,y)$ नाही:
$S(y)$ हा कोणताही एकस्थानी संगणनक्षम संबंध असेल, तर असा काही $k$
असतो की प्रत्येक $y$ साठी $S(y)$ सत्य असते तेव्हाच आणि तरच
$R(k,y)$ सत्य असते.
\end{lem}

\begin{proof}
$R(x,y)$ हा सार्वत्रिक संगणनक्षम संबंध आहे असे समजू.
$S(y)$ हा $\lnot R(y,y)$ हा संबंध घेऊ. $S(y)$ संगणनक्षम असल्याने
काही $k$ साठी $S(y)$ हे $R(k,y)$ शी सममूल्य असेल.
पण मग $S(k)$ हे $R(k,k)$ आणि $\lnot R(k,k)$ या
दोन्हींशी सममूल्य ठरेल; हा व्याघात आहे.
\end{proof}


\begin{thm}
$\Th{T}$ ही $\Th{Q}$ समाविष्ट असलेली कोणतीही सुसंगत उपपत्ती
असेल, तर $\Th{T}$ अनिर्णेय असते.
\end{thm}


\begin{proof}
$\Th{T}$ हा $\Th{Q}$ चा सुसंगत, निर्णेय विस्तार आहे असे समजू.
$\Th{T}$ वापरून सार्वत्रिक संगणनक्षम संबंध परिभाषित करू
आणि त्यामुळे व्याघात मिळवू.

$R(x,y)$ पुढील अटीवर, आणि फक्त त्याच अटीवर, सत्य असू द्या:
\begin{quote}
$x$ हा $!D(u)$ या सूत्राचा संकेतांक आहे, आणि $\Th{T}$ हे
$!D(\num y)$ सिद्ध करते.
\end{quote}
$\Th{T}$ निर्णेय आहे असे गृहीत धरले असल्यामुळे $R$ संगणनक्षम
आहे: सूत्र-संकेतांक ओळखणे व त्यात संख्यांकाचे प्रतिस्थापन करणेही
संगणनक्षम आहे. आता $R$ सार्वत्रिक असल्याचे दाखवू.
$S(y)$ हा कोणताही संगणनक्षम संबंध असेल, तर त्याचे $\Th{Q}$ मध्ये
(आणि म्हणून $\Th{T}$ मध्येही) $!D_S(u)$ या सूत्राने निरूपण होते.
म्हणून प्रत्येक $n$ साठी
\begin{eqnarray*}
S(n) & \lif & T \vdash !D_S(\num n) \\
& \lif & R(\Gn{!D_S(u)}, n)
\end{eqnarray*}
आणि
\begin{eqnarray*}
\lnot S(n) & \lif & T \vdash \lnot !D_S(\num n) \\
& \lif & T \not\vdash !D_S(\num n) \quad \text{($\Th{T}$ सुसंगत असल्यामुळे)} \\
& \lif & \lnot R(\Gn{!D_S(u)}, n).
\end{eqnarray*}
म्हणजे प्रत्येक $y$ साठी $S(y)$ सत्य असते तेव्हाच आणि तरच
$R(\Gn{!D_S(u)}, y)$ सत्य असते. त्यामुळे $R$ सार्वत्रिक
ठरतो आणि अपेक्षित व्याघात मिळतो.
\end{proof}

``सत्य अंकगणित'' ही $\Setabs{!A}{ \Sat{\Nat}{!A}}$ उपपत्ती
असू द्या; म्हणजे अंकगणिताच्या भाषेत प्रमाण अर्थनिर्धारणानुसार
सत्य असलेल्या वाक्यांचा हा संच आहे.
प्रमाण प्रतिमानात क्यूची स्वयंसिद्धके सत्य असल्याने
सत्य अंकगणित हा त्याचा सुसंगत विस्तार आहे.

\begin{cor}
सत्य अंकगणित अनिर्णेय आहे.
\end{cor}

%विभाग~\olref{undefinability:truth:section} मध्ये टार्स्की यांचा
%याहून अधिक सबल निकाल मांडू.

\end{document}
